% =============================================================
%  Chapter 8 — Lossy Reception and Oral Tradition
%  Lean source: FormalAnthropology/IDT_Signal_Ch08.lean
% =============================================================
\chapter{Lossy Reception and Oral Tradition}
\label{ch:lossy}

\epigraph{%
  ``The story becomes at each reproduction more concise, more schematized
  \ldots\ Proper names drop out. Vivid details vanish.
  What remains is a general form.''}%
  {Frederic Bartlett, \textit{Remembering} (1932)}

% ----------------------------------------------------------------
\section{Introduction: The Thermodynamics of Storytelling}
\label{sec:lossy-intro}

In 1932, Frederic Bartlett conducted a series of serial reproduction
experiments that became foundational for cognitive anthropology.  He asked a
subject to read a Native American folk tale (``The War of the Ghosts''), then
reproduce it from memory; the next subject read only the reproduction, and
so on through a chain of ten or more retellings.  The results were dramatic:
details dropped out, unfamiliar concepts were rationalized into familiar
categories, and the story simplified with each retelling until it converged
on a stable, culturally assimilated ``skeleton.''\footnote{Bartlett, F.~C.,
\textit{Remembering: A Study in Experimental and Social Psychology}
(Cambridge, 1932), ch.~5.  The experiment has been replicated hundreds of
times with consistent results; see Mesoudi \& Whiten (2008) for a
meta-analysis.}

Oral tradition---the transmission of myths, genealogies, ritual
instructions, and folk knowledge from person to person without writing---is
the primordial form of cultural transmission.  Before the invention of
writing (c.~3200 BCE in Sumer), \emph{all} culture was oral.  Even today,
the vast majority of human cultural knowledge is transmitted orally: recipes,
gossip, jokes, parenting advice, religious instruction.

This chapter formalizes the defining feature of oral transmission:
\emph{lossiness}.  A receiver who processes a signal through an oral
channel---through memory, paraphrase, reinterpretation---can never produce
output more complex than the input.  Depth can only decrease or remain
constant; it cannot increase.  This is the ``second law of cultural
thermodynamics.''\footnote{The analogy to thermodynamic entropy is
imperfect but illuminating.  In thermodynamics, the second law states that
entropy never decreases in an isolated system.  Here, the analog is that
\emph{complexity} (depth) never increases through lossy reception.  The
``entropy'' of cultural transmission is the \emph{loss} of depth.}

We define three levels of receiver: a generic \emph{Receiver} (preserving
void and distributing over composition), a \emph{LossyReceiver} (depth
non-increasing), and a \emph{StrictlyLossy} receiver (depth strictly
decreasing for complex ideas).  We prove that lossy reception composes,
iterates, and ultimately converges: only shallow ideas survive the
``telephone game.''


% ================================================================
\section{Core Structures}
\label{sec:lossy-defs}

\begin{definition}[Receiver]
\label{def:receiver}
A \emph{receiver} is an endomorphism $\varphi:\IS\to\IS$ satisfying:
\begin{enumerate}
  \item \emph{Void preservation:} $\varphi(\void) = \void$.
  \item \emph{Distributivity:} $\varphi(a\compose b) = \varphi(a)\compose\varphi(b)$.
\end{enumerate}
\end{definition}

A receiver is a ``cultural filter''---a formalized version of how a
particular mind processes incoming signals.  The two axioms are natural:
silence in yields silence out, and hearing a compound message is the same
as hearing the parts and composing the results.

\begin{lstlisting}[caption={The Receiver structure in Lean~4.}]
structure Receiver (I : Type*) [IdeaticSpace I] where
  process : I -> I
  process_void :
    process IdeaticSpace.void = IdeaticSpace.void
  process_compose : forall (a b : I),
    process (IdeaticSpace.compose a b) =
    IdeaticSpace.compose (process a) (process b)
\end{lstlisting}

\begin{definition}[Lossy Receiver]
\label{def:lossy-receiver}
A \emph{lossy receiver} is a receiver $\varphi$ such that
\[
  \depth(\varphi(a)) \;\le\; \depth(a) \quad\text{for all } a\in\IS.
\]
Processing never increases depth.
\end{definition}

This is the defining property of oral transmission: you can never come out
of a retelling with \emph{more} complexity than you went in with.  Written
transmission is different (you can re-read, annotate, cross-reference), but
oral transmission is inherently lossy.

\begin{definition}[Strictly Lossy Receiver]
\label{def:strictly-lossy}
A \emph{strictly lossy receiver} is a lossy receiver $\varphi$ such that
\[
  \depth(a) > 1 \;\implies\;
  \depth(\varphi(a)) < \depth(a).
\]
For complex ideas ($\depth > 1$), depth \emph{strictly} decreases.  Ideas of
depth $\le 1$ may be preserved.
\end{definition}

This models the ``telephone game'': every retelling loses something, unless
the message is already maximally simple.  Only atomic ideas survive
perfect-fidelity transmission through a strictly lossy channel.

\begin{lstlisting}[caption={Strictly lossy receiver in Lean~4.}]
structure StrictlyLossy (I : Type*) [IdeaticSpace I]
    extends LossyReceiver I where
  strict_decay : forall (a : I),
    IdeaticSpace.depth a > 1 ->
    IdeaticSpace.depth
      (toLossyReceiver.toReceiver.process a) <
      IdeaticSpace.depth a
\end{lstlisting}


% ================================================================
\section{Basic Properties of Lossy Reception}
\label{sec:lossy-basic}

\begin{theorem}[Lossy Void Preservation (8.1)]
\label{thm:lossy-void-preservation}
All receivers preserve silence: $\varphi(\void) = \void$.
\end{theorem}

No transmission channel creates content from nothing.  If you say nothing,
the listener hears nothing.  This is the zeroth law of communication---prior
to any thermodynamic considerations, it establishes the baseline: the null
signal is universally transparent.

\begin{theorem}[Lossy Depth Bound (8.2)]
\label{thm:lossy-depth-bound}
After lossy processing: $\depth(\varphi(a)) \le \depth(a)$.
\end{theorem}

An oral retelling is never more complex than the original.  This is the
second law of cultural thermodynamics: complexity never increases through
oral transmission.  The theorem is definitional---it \emph{is} the defining
property of a lossy receiver---but stating it as a theorem emphasizes its
central role.%
\footnote{Contrast with \emph{written} transmission, which can \emph{increase}
depth through annotation, commentary, and cross-reference.  The Talmud is
deeper than the Torah; the scholastic summa is deeper than any single
patristic text.  Writing breaks the second law of oral thermodynamics.}

\begin{theorem}[Lossy Composition Bound (8.3)]
\label{thm:lossy-compose-bound}
Processing a composed idea through a lossy receiver:
\[
  \depth(\varphi(a\compose b)) \;\le\; \depth(a) + \depth(b).
\]
\end{theorem}

\begin{proof}[Proof sketch]
By lossiness, $\depth(\varphi(a\compose b))\le\depth(a\compose b)$.  By
subadditivity of composition depth, $\depth(a\compose b)\le\depth(a)+\depth(b)$.
Chaining the inequalities yields the result.
\end{proof}

A story about two events, retold through oral tradition, has depth at most
the sum of the two events' original depths---even before accounting for the
receiver's lossiness.

\begin{theorem}[Lossy Receiver Preserves Void Depth (8.4)]
\label{thm:lossy-void-depth}
$\depth(\varphi(\void)) = 0.$
\end{theorem}

Silence, when transmitted lossily, remains at depth zero.  Combined with
Theorem~\ref{thm:lossy-void-preservation}, this says not just that void
maps to void, but that the image has zero complexity---a stronger
quantitative statement.


% ================================================================
\section{Iterated Lossy Reception}
\label{sec:lossy-iterated}

What happens when a signal passes through a \emph{chain} of lossy receivers?
Bartlett's serial reproduction experiments provide the empirical motivation;
this section provides the formal theory.

\begin{definition}[Iterated Process]
\label{def:iterate-process}
The $n$-fold iteration of a receiver's process function:
\[
  \varphi^{(0)}(a) = a, \qquad
  \varphi^{(n+1)}(a) = \varphi(\varphi^{(n)}(a)).
\]
\end{definition}

\begin{lstlisting}[caption={Iterated process in Lean~4.}]
def iterateProcess (R : Receiver I) : ℕ -> I -> I
  | 0 => fun a => a
  | n + 1 => fun a => R.process (iterateProcess R n a)
\end{lstlisting}

\begin{theorem}[Iterated Lossy Depth Bound (8.5)]
\label{thm:iterated-lossy-depth}
After $n$ applications of a lossy receiver:
$\depth(\varphi^{(n)}(a)) \le \depth(a)$ for all $n$.
\end{theorem}

\begin{proof}[Proof sketch]
By induction on $n$.  The base case ($n=0$) is trivial.  For the inductive
step, apply the lossiness bound to the result of the $(n)$-th iteration,
whose depth is $\le\depth(a)$ by the inductive hypothesis.
\end{proof}

Oral tradition through a chain of $n$ retellings can only lose depth, never
gain it.  Whether the story passes through 3~or 300~people, it can only get
simpler.  This is the formal backbone of Bartlett's experimental finding.

\begin{theorem}[Iterated Void Stability (8.6)]
\label{thm:iterated-void-stable}
$\varphi^{(n)}(\void) = \void$ for all $n$.
\end{theorem}

\begin{proof}[Proof sketch]
Induction on $n$: the base case is identity, and each step applies
$\varphi(\void)=\void$.
\end{proof}

A chain of listeners hearing nothing produces nothing, no matter how long
the chain.  Silence is the unique \emph{absorbing state} of every receiver.

\begin{theorem}[Lossy Reception of Compositions (8.7)]
\label{thm:lossy-distribute}
$\varphi(a\compose b) = \varphi(a)\compose\varphi(b).$
\end{theorem}

Cultural filtering is ``linear'' in the algebraic sense: a receiver's
reaction to a compound stimulus is the compound of their reactions to the
parts.  This is the distributivity axiom of the Receiver structure, stated
as a theorem for emphasis.

In anthropological terms, it means that a listener's interpretation of a
two-part story can be decomposed into their interpretation of part~1 composed
with their interpretation of part~2.  There is no holistic ``Gestalt'' effect
in lossy reception---only part-by-part processing.%
\footnote{This is a strong structural assumption.  Real human cognition
arguably does exhibit Gestalt effects.  The distributivity axiom is a
simplification that enables clean formal results at the cost of ignoring
non-decomposable cognitive processes.}


% ================================================================
\section{Strict Lossiness and Convergence}
\label{sec:lossy-strict}

\begin{theorem}[Strict Decay Bound (8.8)]
\label{thm:strict-decay}
For a strictly lossy receiver and any idea with $\depth(a)>1$:
\[
  \depth(\varphi(a)) < \depth(a).
\]
\end{theorem}

In ``telephone game'' transmission, every complex message becomes strictly
simpler.  This is not a bug---it's a feature.  Cultural evolution works by
simplification.  The myths that survive thousands of years of oral
transmission are the ones that have been \emph{simplified} to the point
where no further simplification occurs.

\begin{theorem}[Shallow Ideas Survive Lossy Reception (8.9)]
\label{thm:shallow-survives}
If $\depth(a)\le 1$, then $\depth(\varphi(a))\le 1$.
\end{theorem}

\begin{proof}[Proof sketch]
By lossiness: $\depth(\varphi(a))\le\depth(a)\le 1$.
\end{proof}

Simple ideas---basic emotions, universal gestures, atomic symbols---survive
any amount of lossy transmission.  This is why certain cultural universals
persist: they're too simple to be degraded further.  The smile, the frown,
the pointing gesture, the head-nod---all have depth $\le 1$ and are thus
impervious to oral degradation.%
\footnote{Ekman's (1971) cross-cultural studies of facial expression provide
empirical support: the six ``basic emotions'' are recognized across cultures,
arguably because they are too simple to be distorted by oral transmission.}

\begin{theorem}[Zero-Depth Ideas Survive Lossy Reception (8.10)]
\label{thm:zero-depth-survives}
If $\depth(a) = 0$, then $\depth(\varphi(a)) = 0$.
\end{theorem}

Void-like ideas (pure silence, empty gestures) are immune to lossy
transmission.  They are fixed points of every lossy receiver.


% ================================================================
\section{Composition of Receivers}
\label{sec:lossy-composition}

\begin{definition}[Composed Receivers]
\label{def:compose-receivers}
Given receivers $\varphi_1,\varphi_2$, their composition is
$(\varphi_2\circ\varphi_1)(a) = \varphi_2(\varphi_1(a))$---the result of
sequential reception.
\end{definition}

\begin{lstlisting}[caption={Receiver composition in Lean~4.}]
def composeReceivers (R_1 R_2 : Receiver I) : Receiver I where
  process := fun a => R_2.process (R_1.process a)
  process_void := by
    rw [R_1.process_void, R_2.process_void]
  process_compose := by
    intro a b
    rw [R_1.process_compose, R_2.process_compose]
\end{lstlisting}

\begin{theorem}[Composition of Lossy Receivers is Lossy (8.11)]
\label{thm:compose-lossy}
If $\varphi_1$ and $\varphi_2$ are lossy, then
$\varphi_2\circ\varphi_1$ is lossy:
\[
  \depth((\varphi_2\circ\varphi_1)(a)) \;\le\; \depth(a)
  \quad\text{for all } a.
\]
\end{theorem}

\begin{proof}[Proof sketch]
$\depth(\varphi_2(\varphi_1(a)))\le\depth(\varphi_1(a))\le\depth(a)$
by applying lossiness twice.
\end{proof}

A chain of imperfect oral transmitters is itself an imperfect transmitter.
Lossiness is closed under composition---there's no way to ``recover'' lost
depth by adding more lossy steps.  This is the impossibility of error
correction in purely oral tradition.  Unlike digital communication (where
redundancy enables error correction), oral channels have no mechanism for
restoring lost depth.

\begin{theorem}[Double Lossy is More Lossy (8.12)]
\label{thm:double-lossy}
$\depth((\varphi_2\circ\varphi_1)(a)) \le \depth(\varphi_1(a)).$
\end{theorem}

Two generations of oral transmission is worse (or equal) than one.  The
second lossy step can only further reduce what the first step already
degraded.  This is the formal version of the folk wisdom that stories ``get
worse with each retelling.''


% ================================================================
\section{Resonance and Lossy Reception}
\label{sec:lossy-resonance}

\begin{theorem}[Receiver Preserves Self-Resonance (8.13)]
\label{thm:receiver-self-resonance}
$\varphi(a) \res \varphi(a)$ for all $a$.
\end{theorem}

Even lossy transmission produces internally consistent output.  A garbled
retelling is garbled, but it still ``makes sense to itself.''  This is
a consequence of the reflexivity of resonance, but it carries anthropological
weight: oral traditions, however degraded, are never \emph{incoherent}.
They simplify, regularize, and schematize, but they always remain
self-consistent.

\begin{theorem}[Lossy Compound Self-Resonance (8.14)]
\label{thm:lossy-compound-resonance}
$\varphi(a\compose b) \res \varphi(a\compose b).$
\end{theorem}

Even through lossy reception, the receiver's output from a compound signal
is internally coherent.  The theorem is trivial (self-resonance is always
true), but it anchors the stronger results of Chapter~\ref{ch:multireceiver},
where we study how \emph{different} receivers' outputs relate to each other.

\begin{theorem}[Void Iterate Always Resonates (8.15)]
\label{thm:void-iterate-resonates}
$\varphi^{(n)}(\void) \res \varphi^{(n)}(\void)$ for all $n$.
\end{theorem}

Chains of silence produce silence, which trivially resonates with itself.
Combined with Theorem~\ref{thm:iterated-void-stable}, this says that
void is not only a fixed point of iteration but a self-resonant one.


% ================================================================
\section{Advanced Properties}
\label{sec:lossy-advanced}

\begin{theorem}[Lossy Iterated Composition Bound (8.16)]
\label{thm:lossy-composeN}
$\depth(\varphi(\compN{a}{n})) \le n\cdot\depth(a).$
\end{theorem}

\begin{proof}[Proof sketch]
By lossiness, $\depth(\varphi(\compN{a}{n}))\le\depth(\compN{a}{n})$.
By the composeN depth bound, $\depth(\compN{a}{n})\le n\cdot\depth(a)$.
\end{proof}

A repeated message (ritual chant, liturgical refrain) loses depth through
lossy reception, bounded by the repetition count times original depth.  Even
before lossiness takes its toll, the starting depth is only linearly bounded
in repetition count.

\begin{theorem}[Receiver Composition is Associative on Output (8.17)]
\label{thm:receiver-assoc}
$\varphi_3(\varphi_2(\varphi_1(a)))
= (\varphi_3\circ\varphi_2\circ\varphi_1)(a).$
\end{theorem}

A three-generation transmission chain (grandparent $\to$ parent $\to$ child)
produces the same result regardless of how we ``bracket'' the generations.
This is the associativity of function composition, applied to the receiver
setting.  It ensures that our model of multi-generational oral tradition is
well-defined: the order of composition doesn't depend on arbitrary grouping.

\begin{theorem}[Depth Monotonicity Through Chains (8.18)]
\label{thm:chain-depth-monotone}
For $m\le n$:
$\depth(\varphi^{(n)}(a)) \le \depth(\varphi^{(m)}(a)).$
\end{theorem}

\begin{proof}[Proof sketch]
By induction on $n$.  If $m=n$, the result is trivial.  If $m<n$, apply
lossiness at step $n$ to reduce to the case $n-1$, then apply the
inductive hypothesis.
\end{proof}

At every point in a chain of oral transmission, the story is no more complex
than it was at any earlier point.  There is no ``regeneration'' through
purely lossy channels.  The depth sequence $\depth(\varphi^{(0)}(a)),
\depth(\varphi^{(1)}(a)), \depth(\varphi^{(2)}(a)), \ldots$ is monotonically
non-increasing.

This is the formal version of a deep anthropological observation: myths
don't get \emph{worse}; they get \emph{stable}.  Oral tradition doesn't
``degrade''---it \emph{converges} to a fixed point determined by the
transmission medium's lossiness.  The stories that survive are the ones
simple enough to be transmitted without further loss.

\begin{lstlisting}[caption={Chain depth monotonicity in Lean~4.}]
theorem chain_depth_monotone
    (R : LossyReceiver I) (a : I) (m n : ℕ)
    (hmn : m <= n) :
    IdeaticSpace.depth
      (iterateProcess R.toReceiver n a) <=
    IdeaticSpace.depth
      (iterateProcess R.toReceiver m a)
\end{lstlisting}


% ================================================================
\section{Conclusion}
\label{sec:lossy-conclusion}

This chapter established the formal theory of lossy reception, providing
mathematical foundations for Bartlett's serial reproduction experiments and
the broader study of oral tradition.

The central results are:
\begin{enumerate}
  \item Lossy reception is closed under composition
    (Theorem~\ref{thm:compose-lossy}): chains of oral transmitters are
    themselves oral transmitters.
  \item Depth is monotonically non-increasing through chains
    (Theorem~\ref{thm:chain-depth-monotone}): stories only get simpler.
  \item Shallow ideas survive (Theorem~\ref{thm:shallow-survives}):
    cultural universals persist because they are too simple to degrade.
\end{enumerate}

The theory explains why oral traditions converge: not because storytellers
are ``bad'' at their job, but because lossiness is a structural feature of
the oral channel.  Written traditions escape this constraint by providing
external memory, but oral traditions are bound by the second law of cultural
thermodynamics.

Chapter~\ref{ch:echo} takes this analysis inward: what happens when a
community applies lossy reception \emph{to its own output}, creating an echo
chamber?
